\subsubsection{Répartition des fronts d’onde entre accélérateur et processeur}\label{sec:wfa-gpu-batch}
WFA-GPU, développé par Aguado-Puig et coll.~\cite{wfagpuaffine,wfagpucode}, distribue les alignements affines entre blocs CUDA et les diagonales d’un front entre fils. Soient \(x\) le coût d’une substitution, \(o\) celui de l’ouverture d’une lacune et \(e\) celui de son extension. Pour un chemin \(P\), le coût dépend de l’ensemble des lacunes \(G(P)\) et de leurs longueurs \(\ell_g\),
\[
 C(P)=xN_{\mathrm{sub}}(P)+\sum_{g\in G(P)}(o+e\ell_g),\qquad
 (x,o,e)=(1,0,1)\ \Longrightarrow\ \min_P C(P)=d(A,B).
\]
Le mode score utilise le front complet. Le GPU traite au plus \(512\) paires par lot. Lorsqu’un front atteint la limite de calcul ou qu’une chaîne contient \texttt{N}, WFA-GPU confie la paire à son sous-module WFA2~\cite{wfa2gpucode}. Le symbole \texttt{N} conserve sa valeur littérale.

La figure~\ref{fig:wfagpuaffine} compare cinq blocs après un bloc de chauffe sur GH200. Chaque temps couvre le lancement du processus, la lecture des chaînes et l’écriture des scores. À \(20\)~knt et \(10\,\%\) d’éditions injectées, le rapport des temps WFA-GPU/WFA2 passe de \(15{,}77\) pour cinq paires à \(0{,}762\) pour les \(128\) paires qui les contiennent. Sur ce dernier lot, WFA-GPU prend \(0{,}395\pm0{,}006\)~s, edlib \(0{,}250\pm0{,}001\)~s et NvPD CUDA \(2{,}524\pm0{,}028\)~s. NvPD parcourt le rectangle de chaque paire, tandis que la largeur du front dépend du score atteint.
\chart{gpu/lot_fronts_affines.png}{Rapport des temps de processus WFA-GPU et WFA2 sur un cœur Grace pour les lots synthétiques. Les points donnent les moyennes des rapports appariés, les traits épais l’écart-type et les traits fins l’intervalle de Student à 95\,\%. Le panneau b répartit les paires biologiques entre CPU et GPU. Les effectifs sont identiques dans les cinq blocs.}{fig:wfagpuaffine}{python/studies/wfa_gpu_batch.py}
\FloatBarrier
Sur les \(2\,016\) paires de \(\phi\)X174, les trois génomes contenant \texttt{N} participent à \(186\) paires. Le programme les reprend toutes sur CPU. Sur le chromosome~21, les \(768\) fenêtres homologues de \(16\,384\) bases de référence ont des distances comprises entre \(0\) et \(11\,107\). Porter la limite \(K\) de \(64\) à \(1\,024\) fait passer le nombre de reprises de \(152\) à \(10\). Le rapport apparié \(T_{1024}/T_{64}=0{,}958\) a pour intervalle à 95\,\% \([0{,}899\,;\,1{,}017]\). Le réglage change aussi la taille de bloc que le programme choisit pour le GPU.

Pour la matrice SARS-CoV-2, le pilote réserve le GPU aux paires dont chaque chaîne compte moins de \(32\,768\) bases et appartient à \(\{A,C,G,T\}^{*}\). Les six génomes admis fournissent \(\binom{6}{2}=15\) paires sur \(2\,016\). Quatre travailleurs CPU utilisant WFA2-lib 2.3.6~\cite{wfa2code} calculent les \(2\,001\) autres paires avec les symboles IUPAC littéraux, pendant que le GPU calcule les quinze scores. Les témoins CPU utilisent les mêmes quatre travailleurs. Du FASTA au TSV, le calcul hybride prend \(6{,}844\pm0{,}015\)~s, WFA2 \(6{,}874\pm0{,}030\)~s et edlib \(2{,}628\pm0{,}014\)~s. L’intervalle du rapport hybride/WFA2 vaut \([0{,}9908\,;\,1{,}0004]\).

À coût unitaire, le front de score \(s\) admet au plus \(2s+1\) diagonales, avant troncature aux extrémités des chaînes. Pour une paire de distance \(d\), le nombre cumulé de positions possibles vaut donc \(H(d)=\sum_{s=0}^{d}(2s+1)=(d+1)^2\). Les quinze paires admissibles sur GPU ne représentent que \(223\,071\) des \(26\,084\,927\,678\) positions de ce modèle. Celui-ci omet les extensions de correspondances. Dans la mesure, le lot GPU prend en moyenne \(0{,}498\)~s et le lot CPU le plus lent \(6{,}684\)~s.

\synthese{Sur le lot synthétique le plus divergent à \(128\) paires de \(20\)~knt, le rapport à WFA2 a pour intervalle à 95\,\% \([0{,}748\,;\,0{,}776]\). Sur les fenêtres chromosomiques, le meilleur temps moyen de WFA-GPU reste \(14{,}12\) fois celui de WFA2 sur un cœur Grace.}
